\chapter{Generalised Linear Models and the Regression Family}
\companion{StatQuest: Generalised Linear Models (search)}{\yts{StatQuest+Generalized+Linear+Models}}

The InfoEdge hire in the experience video said the single biggest change in his understanding came from
learning the ``variants most courses skip'': count data needs \emph{Poisson regression}, which comes from
\emph{GLMs}; SVMs can do \emph{regression}. In Round 4 he answered a heteroscedasticity question with a GLM and the
interviewer accepted it. This chapter makes you fluent in that family. It is a high-return chapter because most
candidates have never seen it.

\section{The question GLMs answer}

Ordinary linear regression quietly assumes three things about the target $y$:
\begin{enumerate}
\item its average is a straight-line function of the features;
\item it can be any real number, negative included;
\item its noise has the same spread everywhere (and is roughly bell-shaped).
\end{enumerate}
Many real targets break all three:
\begin{itemize}
\item \textbf{Counts}: number of applications a job gets today (0, 1, 2, \dots). Never negative; a job that
  averages 2 applies varies by about $\pm1.4$, one that averages 200 varies by about $\pm14$. Spread grows with the mean.
\item \textbf{Positive, skewed amounts}: salary, property price, time-to-hire. Never negative; spread grows in
  proportion to the level.
\item \textbf{Yes/no and proportions}: did the user apply; what fraction of mails were opened. Must stay in $[0,1]$.
\end{itemize}
A \textbf{generalised linear model (GLM)} keeps the convenient ``weighted sum of features'' but lets us choose
(a) what distribution the target follows and (b) how the weighted sum is connected to the target's average.

\section{The three ingredients of a GLM}

\begin{enumerate}
\item \textbf{A linear predictor} (the familiar part): $\eta = \mathbf w^\top\mathbf x + b$. It can be any real number.
\item \textbf{A link function} $g$ that connects the average of $y$, written $\mu = \E[y\mid\mathbf x]$, to the linear
  predictor: $g(\mu) = \eta$, or equivalently $\mu = g^{-1}(\eta)$. The link's job is to map the allowed range of $\mu$
  onto the whole real line.
\item \textbf{A distribution} for $y$ around its mean, from the \emph{exponential family} (Gaussian, Bernoulli/Binomial,
  Poisson, Gamma, Inverse Gaussian, Tweedie, \dots). The distribution fixes how the variance depends on the mean:
  $\Var(y) = \phi\,V(\mu)$, where $V$ is the \textbf{variance function} and $\phi$ a dispersion constant.
\end{enumerate}

\begin{center}\small
\begin{tabularx}{\linewidth}{l l l l X}
\toprule
Target type & Distribution & Variance $V(\mu)$ & Usual link & Name\\
\midrule
Any real number & Gaussian & $1$ (constant) & identity $\mu = \eta$ & ordinary linear regression\\
0/1 or proportion & Bernoulli / Binomial & $\mu(1-\mu)$ & logit $\log\frac{\mu}{1-\mu} = \eta$ & logistic regression\\
Count $0,1,2,\dots$ & Poisson & $\mu$ & log $\log\mu = \eta$ & Poisson regression\\
Overdispersed count & Negative binomial & $\mu + \mu^2/k$ & log & NB regression\\
Positive, skewed & Gamma & $\mu^2$ & log (or inverse) & Gamma regression\\
Zeros + positive amounts & Tweedie ($1<p<2$) & $\mu^p$ & log & Tweedie regression\\
\bottomrule
\end{tabularx}
\end{center}

\begin{intuition}[: the link is a translator, the distribution is a noise model]
The linear predictor speaks ``any real number''. A count's mean speaks ``positive numbers''. The log link translates:
$\mu = e^\eta$ is always positive, whatever $\eta$ is. The distribution then says how scattered the real counts are
around that mean. Ordinary regression is the special case where no translation is needed and the scatter is the same
everywhere.
\end{intuition}

\section{Poisson regression, built from scratch}

\subsection{The Poisson distribution in one paragraph}
If events happen independently at an average rate $\mu$ per interval, the count $Y$ in one interval follows
$P(Y = k) = e^{-\mu}\mu^k/k!$ for $k = 0,1,2,\dots$ Its mean \emph{and} its variance are both $\mu$.
\begin{example}[: Poisson probabilities]
A job averages $\mu = 2$ applies per hour. $P(0) = e^{-2} = 0.135$; $P(1) = 2e^{-2} = 0.271$; $P(2) = 2e^{-2} = 0.271$;
$P(3) = \frac{8}{6}e^{-2} = 0.180$.
\end{example}

\subsection{The model}
\[ y\mid\mathbf x \sim \text{Poisson}(\mu),\qquad \log\mu = \mathbf w^\top\mathbf x + b\quad\Longleftrightarrow\quad \mu = e^{\mathbf w^\top\mathbf x + b}. \]
Effects are \emph{multiplicative}: raising $x_j$ by one unit multiplies the expected count by $e^{w_j}$.

\begin{example}[: reading a Poisson coefficient]
Daily applies to a job: $\log\mu = 1.5 + 0.4\cdot\text{remote} - 0.05\cdot\text{days\_since\_posting}$.
A fresh, non-remote job: $\mu = e^{1.5} = 4.48$ applies/day. Remote: $e^{0.4} = 1.49$, so 49\% more, $6.7$/day.
After 10 days: multiply by $e^{-0.5} = 0.61$.
\end{example}

\subsection{How it is fitted}
Maximum likelihood again. For one row, $\log P(y\mid\mu) = y\log\mu - \mu - \log y!$. With $\log\mu = \eta$, the
negative log-likelihood to minimise (dropping the constant) is
\[ L(\mathbf w) = \sum_i\big(e^{\eta_i} - y_i\eta_i\big). \]
Its gradient with respect to $\mathbf w$ is $\sum_i(\mu_i - y_i)\mathbf x_i$: \emph{prediction minus truth, times input},
exactly the same form as linear and logistic regression. This is no coincidence: it is a property of every GLM that uses
its \textbf{canonical link}, and it is why one algorithm (IRLS / Newton) fits them all. The loss is convex, so the optimum
is global.

\subsection{Exposure and offsets}
Counts often come from different-sized windows: one job was live for 3 days, another for 30. Model the \emph{rate}:
$\log\mu_i = \log(\text{exposure}_i) + \mathbf w^\top\mathbf x_i$. The term $\log(\text{exposure})$ is an \textbf{offset}: a
feature whose coefficient is fixed at 1. Now $e^{w_j}$ is a rate ratio.

\subsection{Overdispersion and zero inflation}
Real counts usually vary \emph{more} than Poisson allows (variance $>$ mean), because of unobserved differences between
units. Symptoms: residual deviance much larger than its degrees of freedom; standard errors far too small.
\begin{itemize}
\item \textbf{Quasi-Poisson}: keep the Poisson mean model, scale the variance by $\phi > 1$ (fixes standard errors).
\item \textbf{Negative binomial}: variance $\mu + \mu^2/k$; a Poisson whose rate itself varies (Gamma-distributed).
\item \textbf{Zero-inflated} or \textbf{hurdle} models: when there are far more zeros than Poisson predicts (most jobs
  get zero applies from a given city), model ``zero or not'' and ``how many, if any'' separately.
\end{itemize}

\section{Gamma regression: positive, skewed, multiplicative noise}
The Gamma distribution lives on positive numbers and has variance proportional to $\mu^2$, i.e. a \emph{constant
coefficient of variation}: a 10 lakh salary varies by about $\pm1.5$ lakh and a 100 lakh salary by about $\pm15$
lakh. With a log link, $\log\E[y] = \mathbf w^\top\mathbf x + b$.

Why prefer it over ``regress $\log y$ with OLS''?
\begin{itemize}
\item OLS on $\log y$ estimates $\E[\log y]$, and $e^{\E[\log y]}$ is the geometric mean, which is \emph{below} the
  arithmetic mean $\E[y]$. You need a correction to predict average salary. Gamma-log models $\E[y]$ directly.
\item It handles the heteroscedasticity through the variance function rather than through a transformation.
\end{itemize}
This is exactly the answer the candidate gave for violated homoscedasticity: ``model the variance explicitly, keep a
linear model on top''.

\section{Tweedie: when many values are exactly zero}
Revenue per recruiter this month, insurance claims, money spent by a user: most are exactly 0, the rest positive and
skewed. The Tweedie family with power $1<p<2$ is a compound Poisson--Gamma: a Poisson number of Gamma-sized amounts, so
it has a point mass at 0 and a continuous positive part. LightGBM and XGBoost offer \code{objective="tweedie"}.

\section{Beyond the mean: other regression losses}

GLMs change the distribution. Another way to vary regression is to change \emph{what summary of $y$} you predict, or
\emph{how you treat large errors}.

\subsection{Quantile regression}
Instead of the mean, predict a quantile: the 90th percentile salary for a role, so that ``90\% of offers are below this''.
The \textbf{pinball (quantile) loss} for quantile $\tau$ is
\[ \ell_\tau(y,\hat y) = \begin{cases}\tau\,(y - \hat y) & y \ge \hat y\\ (1-\tau)(\hat y - y) & y < \hat y\end{cases}. \]
Under-prediction costs $\tau$ per unit and over-prediction $1-\tau$. At $\tau = 0.5$ it is half the absolute error, and
the minimiser is the median.
\begin{example}[: why the pinball loss finds the 90th percentile]
Values $1,2,\dots,10$, $\tau = 0.9$. If $\hat y = 9$: one value above (cost $0.9\times1$) and eight below (cost $0.1\times
(8+7+\dots+1) = 3.6$), total $4.5$. If $\hat y = 10$: nothing above, cost $0.1\times(9+8+\dots+1)=4.5$. If $\hat y = 5$:
above costs $0.9\times(1+2+3+4+5) = 13.5$, below $0.1\times10 = 1$; total $14.5$. The minimum is around the 90th
percentile, where 90\% of points are below.
\end{example}
Uses: \textbf{salary ranges} (10th--90th percentile band in a salary engine), delivery-time promises, prediction intervals
from gradient boosting (\code{loss="quantile"}).

\subsection{Robust regression: Huber and friends}
\textbf{Huber loss} is quadratic for small errors ($|e| \le \delta$) and linear for large ones:
$\ell = \frac12e^2$ if $|e|\le\delta$, else $\delta(|e| - \frac12\delta)$. Small errors are treated like OLS (efficient);
large errors are treated like MAE (outliers cannot dominate). Other robust options: least absolute deviations (median
regression), RANSAC (fit on random subsets, keep the model with most inliers), Theil--Sen.

\subsection{Support Vector Regression (SVR)}
The SVM idea applied to regression. Use the \textbf{$\varepsilon$-insensitive loss}: errors smaller than $\varepsilon$ cost
\emph{nothing}; beyond that, cost grows linearly:
\[ \ell_\varepsilon(y,\hat y) = \max(0, |y - \hat y| - \varepsilon). \]
Picture a tube of half-width $\varepsilon$ around the prediction. Points inside the tube are ignored; points on or outside
it become \textbf{support vectors} and alone determine the fit. Add an L2 penalty $\frac12\norm{\mathbf w}^2$ for flatness:
\[ \min_{\mathbf w,b}\ \tfrac12\norm{\mathbf w}^2 + C\sum_i\max(0,|y_i - \mathbf w^\top\mathbf x_i - b| - \varepsilon). \]
Because the dual form involves only dot products, \textbf{kernels} apply (RBF SVR fits smooth curves). Hyperparameters:
$C$ (penalty for points outside the tube), $\varepsilon$ (tube width, the tolerated error), and the kernel's $\gamma$.
Chapter 9 develops SVMs fully.

\begin{example}[: the $\varepsilon$-insensitive loss]
$\varepsilon = 0.5$. Errors $0.3, -0.5, 1.2, -2.0$ cost $0, 0, 0.7, 1.5$.
\end{example}

\subsection{Other regressors to be able to name}
\begin{itemize}
\item \textbf{Kernel ridge regression}: ridge regression with the kernel trick; squared loss, so every point matters
  (unlike SVR's sparsity); closed-form solution $\boldsymbol\alpha = (K + \lambda I)^{-1}\mathbf y$.
\item \textbf{Gaussian process regression}: a Bayesian kernel method that returns a mean \emph{and} an uncertainty.
\item \textbf{Isotonic regression}: fits the best monotone (non-decreasing) step function; used to calibrate classifier
  probabilities.
\item \textbf{$k$-NN regression, regression trees, random forests, gradient boosting, neural networks}: all can regress.
\item \textbf{Survival (time-to-event) models}: time-to-hire or time-to-churn with censoring (jobs still open), e.g. Cox
  proportional hazards. Ordinary regression on ``days open'' is biased because open jobs have not finished.
\end{itemize}

\section{Choosing a regression model: a decision guide}
\begin{center}\small
\begin{tabularx}{\linewidth}{X X}
\toprule
Situation & Reach for\\
\midrule
Continuous, roughly symmetric noise & OLS / Ridge\\
Counts (applies, clicks, calls) & Poisson (or negative binomial) with log link and exposure offset\\
Positive, right-skewed, spread $\propto$ level (salary, price) & Gamma GLM with log link, or OLS on $\log y$\\
Many exact zeros plus positive amounts (revenue) & Tweedie, or a two-part (hurdle) model\\
Need a range or a pessimistic estimate & Quantile regression (or quantile GBM)\\
Outliers in $y$ & Huber / LAD / RANSAC\\
Non-linear, smallish data, smooth function & SVR or kernel ridge with RBF kernel, or Gaussian process\\
Big tabular data with interactions & Gradient boosting (with the matching objective: poisson, gamma, tweedie, quantile)\\
Time-to-event with censoring & Survival models\\
\bottomrule
\end{tabularx}
\end{center}

\begin{practical}[: GLMs at InfoEdge]
\begin{itemize}
\item \textbf{Applies per job per day} (Job Posting engine; 100 crore+ applies a year): Poisson/negative binomial with
  exposure, used to forecast which jobs will under-perform and need promotion via mailers or WhatsApp.
\item \textbf{Salary engine} (Data Factory): Gamma-log or log-OLS for the mean, quantile regression for the band shown in
  Naukri TopTier's ``salary intelligence'' (10th/50th/90th percentile).
\item \textbf{99acres}: number of enquiries per listing (Poisson); price per square foot (Gamma or log-OLS).
\item \textbf{Revenue per recruiter} (many zero months): Tweedie.
\end{itemize}
\end{practical}

\stuck{StatQuest: Support Vector Machines Part 1 (for SVR intuition)}{\ytid{efR1C6CvhmE}}
\section{Interview questions}

\begin{qa}{\pyq{reported by InfoEdge hire: count data} You need to predict the number of applications a job posting
will receive tomorrow. Why is ordinary linear regression a poor choice, and what would you use?}
\oneline{Counts are non-negative integers whose variance grows with their mean, and effects on them are naturally
multiplicative; Poisson regression (a GLM with a log link), or negative binomial if overdispersed, models all three.}
\why
Problems with OLS on counts:
\begin{enumerate}
\item It can predict negative applies for unattractive jobs.
\item Its noise model assumes the same spread for a job expecting 1 apply and one expecting 500; in reality the spread
  grows like $\sqrt\mu$. OLS over-weights the noisy high-count jobs and its intervals are wrong.
\item Effects are multiplicative in reality (``remote jobs get 50\% more applies'', not ``7 more applies'' regardless
  of baseline). A linear model forces additivity.
\end{enumerate}
Poisson regression: $y\sim\text{Poisson}(\mu)$, $\log\mu = \mathbf w^\top\mathbf x + b + \log(\text{hours live})$ (offset
for exposure). Predictions are positive; $e^{w_j}$ is a rate ratio; the variance equals the mean by assumption.
\textbf{Check overdispersion}: if the Pearson $\chi^2$ divided by degrees of freedom is well above 1 (say 3), switch to
negative binomial or quasi-Poisson. If there are many structural zeros (expired jobs), use a zero-inflated model.
\worked Fitted $\log\mu = 0.8 + 0.4\,\text{remote} + 0.3\,\text{top\_company} - 0.07\,\text{days}$. A remote job at a
top company on day 0: $\mu = e^{1.5} = 4.5$ applies/day; on day 10: $4.5\times e^{-0.7} = 2.2$.
\pushes \emph{``Could you just log-transform and use OLS?''} $\log(0)$ is undefined; $\log(y+1)$ is a hack that biases
results for small counts. \emph{``Could gradient boosting do it?''} Yes, with \code{objective="poisson"}, which uses the
same log link and Poisson deviance loss, capturing non-linearities.
\pitfall Saying ``use classification because counts are discrete''. Counts are ordered and unbounded; a regression with a
count distribution is the right framing.
\end{qa}

\begin{qa}{\pyq{reported: GLM as the answer to heteroscedasticity} What is a generalised linear model? Name its
components and explain how it generalises linear regression.}
\oneline{A GLM has a linear predictor $\eta = \mathbf w^\top\mathbf x + b$, a link function $g$ with $g(\E[y]) = \eta$, and a
distribution from the exponential family that fixes how the variance depends on the mean; linear regression is the
special case with Gaussian noise and identity link.}
\why
Linear regression ties together three choices: the mean is linear, the target is unbounded, the noise is constant and
Gaussian. A GLM unties them. The link handles the range of the mean (log for positive, logit for $(0,1)$). The
distribution handles the noise shape and the mean--variance relationship (Poisson: $\Var = \mu$; Gamma: $\Var \propto
\mu^2$; Bernoulli: $\mu(1-\mu)$). Fitting is by maximum likelihood, done with iteratively reweighted least squares:
each Newton step is a weighted linear regression on a ``working response''. Goodness of fit is measured by
\textbf{deviance}, twice the log-likelihood gap between the model and a perfect (saturated) model; for the Gaussian it
reduces to the sum of squared residuals.
\worked Three GLMs on the same features: Gaussian-identity for ``rating 1--5'' treated as continuous,
Bernoulli-logit for ``applied or not'', Poisson-log for ``number of profile views''.
\pushes \emph{``Why must the distribution be in the exponential family?''} Because then the likelihood has a common form
$\exp\{(y\theta - b(\theta))/\phi + c(y,\phi)\}$, which gives a single fitting algorithm, convexity with the canonical
link, and the neat gradient $(\mu - y)\mathbf x$.
\end{qa}

\begin{qa}{\pyq{reported: SVR} Can an SVM be used for regression? Explain how.}
\oneline{Yes: support vector regression fits a function that keeps as many points as possible inside a tube of half-width
$\varepsilon$, ignores errors smaller than $\varepsilon$, penalises larger errors linearly with weight $C$, keeps the function
flat with an L2 penalty, and uses kernels for non-linear fits.}
\why
Classification SVM: find the widest margin separating classes; only points on or inside the margin (support vectors)
matter. SVR flips the picture: we want a function such that points lie \emph{within} a band around it. Formally
\[ \min\ \tfrac12\norm{\mathbf w}^2 + C\sum_i(\xi_i + \xi_i^*)\ \text{ s.t. } y_i - f(\mathbf x_i) \le \varepsilon + \xi_i,\
f(\mathbf x_i) - y_i \le \varepsilon + \xi_i^*,\ \xi,\xi^*\ge0. \]
$\xi$ and $\xi^*$ measure how far a point sticks out above or below the tube. The dual involves only dot products
$\mathbf x_i^\top\mathbf x_j$, so replacing them with a kernel $K(\mathbf x_i,\mathbf x_j)$ gives non-linear regression.
The solution depends only on points on or outside the tube: \textbf{sparse} in the data, like SVM classification.
\textbf{Hyperparameters:} $\varepsilon$ (how much error is tolerated for free; larger $\varepsilon$ means fewer support vectors and
a smoother fit), $C$ (how strongly violations are penalised), $\gamma$ (RBF width). Scale features first.
\worked $\varepsilon = 1$ lakh for a salary model: predictions within 1 lakh incur no loss, so the model concentrates on the
large errors; points within the tube do not influence the fit.
\pushes \emph{``SVR vs kernel ridge?''} Same kernel trick; kernel ridge uses squared loss (every point matters, closed
form, faster to fit for moderate $n$), SVR uses $\varepsilon$-insensitive loss (sparse, robust to small noise, linear in large
errors so less outlier-sensitive). \emph{``When is SVR a good choice?''} Small to medium data (up to tens of thousands of
rows), smooth non-linear relations, a natural tolerance level.
\pitfall Saying ``SVMs are only for classification'', or ``points inside the tube are support vectors'' (it is the
reverse).
\end{qa}

\begin{qa}{\deep Is logistic regression a GLM? Is linear regression? What is a ``link function'' in each?}
\oneline{Both are GLMs: linear regression is Gaussian with the identity link, logistic regression is Bernoulli with the
logit link $\log\frac{\mu}{1-\mu}$.}
\why For a binary $y$, $\mu = P(y=1)$ lives in $(0,1)$. The logit maps $(0,1)$ onto the real line, so the linear
predictor can be any number. Its inverse is the sigmoid. The Bernoulli variance $\mu(1-\mu)$ is the variance function.
Other links for binary data exist: \textbf{probit} (inverse normal CDF, popular in economics) and \textbf{complementary
log-log} (asymmetric, for rare events). Logit is the canonical link, giving the tidy gradient and interpretable odds ratios.
\end{qa}

\begin{qa}{\deep Why does every canonical-link GLM have the gradient ``(prediction $-$ target) $\times$ input''?}
\oneline{Because in the exponential family the derivative of the log-partition function is the mean, so with the
canonical link the score equation becomes $\sum(y_i - \mu_i)\mathbf x_i = 0$.}
\why An exponential-family density can be written $p(y) = h(y)\exp(\theta y - A(\theta))$, where $\theta$ is the natural
parameter and $A$ the log-partition function, with $A'(\theta) = \mu$. The canonical link sets $\theta = \eta = \mathbf
w^\top\mathbf x$. The log-likelihood of one point is $\theta y - A(\theta)$, whose derivative with respect to $\mathbf w$ is
$(y - A'(\theta))\mathbf x = (y - \mu)\mathbf x$. Gaussian ($A = \theta^2/2$), Bernoulli ($A = \log(1+e^\theta)$) and
Poisson ($A = e^\theta$) all fit this.
\pushes \emph{``Why does that matter?''} One generic fitting routine (IRLS) and one gradient code path serve all of them;
it is also why neural-network output layers pair sigmoid with BCE and softmax with cross-entropy.
\end{qa}

\begin{qa}{\deep Salary is right-skewed and heteroscedastic. Compare: OLS on $\log(\text{salary})$, a Gamma GLM with log link,
and gradient boosting with MSE.}
\oneline{Log-OLS and Gamma-log both model multiplicative effects and growing spread, but log-OLS predicts the geometric
mean and needs a back-transform correction while Gamma-log predicts the mean directly; GBM with MSE fits non-linear effects
but is dominated by the largest salaries unless you change its objective or target.}
\why
\begin{itemize}
\item \textbf{Log-OLS}: fit $\E[\log y]$. Predicting $e^{\hat\eta}$ underestimates the mean by a factor $e^{\sigma^2/2}$
  (Duan smearing fixes it). Coefficients are percent effects. Easy and robust.
\item \textbf{Gamma-log}: $\log\E[y] = \eta$ and $\Var\propto\mu^2$. Coefficients are percent effects on the mean; no
  back-transform; requires $y > 0$.
\item \textbf{GBM with MSE on raw salary}: captures interactions (city $\times$ title) but squared error on raw scale makes
  the high earners dominate the loss. Better: GBM on $\log y$, or GBM with \code{objective="gamma"}.
\end{itemize}
Choose by the question: for an interpretable salary report, log-OLS or Gamma; for maximum accuracy, GBM with the gamma
objective, validated with a scale-aware metric (MAPE or RMSLE).
\end{qa}

\begin{qa}{What is overdispersion and how do you detect and handle it?}
\oneline{Overdispersion is count data varying more than the Poisson model's variance $=$ mean allows; detect it with
deviance or Pearson $\chi^2$ per degree of freedom well above 1, and handle it with quasi-Poisson, negative binomial, or
random effects.}
\why Unmodelled differences (some recruiters write better job descriptions) make counts clump. Under Poisson the
standard errors are then too small, so you would wrongly declare features significant. Negative binomial adds a Gamma-
distributed multiplier to each rate, giving $\Var = \mu + \mu^2/k$.
\worked Mean daily applies 5, sample variance 23: ratio 4.6, strong overdispersion.
\end{qa}

\begin{qa}{\deep You need a salary \emph{range} for a job title, not a single number. How would you build it?}
\oneline{Use quantile regression (linear or gradient-boosted) to predict, for example, the 10th, 50th and 90th
percentiles directly with the pinball loss; alternatively predict mean and variance with a distributional model.}
\why A mean $\pm$ constant band is wrong under heteroscedasticity. Quantile regression lets each quantile depend on features
differently (senior roles have wider bands). Train one model per quantile with loss $\ell_\tau$; enforce non-crossing (sort
the predictions or fit jointly). Evaluate with pinball loss and \textbf{coverage}: on validation data, about 80\% of actual
salaries should fall between the 10th and 90th predictions.
\inpractice This is what a ``salary intelligence'' card (Naukri TopTier: ``you earn more than 68\% of software engineers in
Delhi'') needs: a conditional distribution, not a point.
\end{qa}

\begin{qa}{Explain the pinball loss and prove that, for $\tau = 0.5$, its minimiser is the median.}
\oneline{It charges $\tau$ per unit of under-prediction and $1-\tau$ per unit of over-prediction; at $\tau=0.5$ both are $0.5$,
so it is half the absolute error, which is minimised by the median.}
\why Minimise $\sum_i\ell_\tau(y_i, c)$ over a constant $c$. Moving $c$ up by a tiny $\delta$ changes the loss by $\delta$
times $[(1-\tau)\cdot\#\{y_i < c\} - \tau\cdot\#\{y_i > c\}]$. At the optimum this is zero: the fraction of points below
$c$ equals $\tau$. So $c$ is the $\tau$-quantile; for $\tau = 0.5$, the median.
\end{qa}

\begin{qa}{\deep What is Huber loss and when would you use it instead of MSE or MAE?}
\oneline{Huber is squared for small errors and linear for large ones; use it when the data are mostly well-behaved but have
some outliers, because it keeps OLS's efficiency on the bulk and MAE's robustness on the tail.}
\why MSE's gradient is the error itself, so one huge error dominates the update. MAE's gradient is $\pm1$ everywhere: robust
but not smooth at 0 and less efficient for Gaussian-like noise. Huber's gradient is the error, clipped at $\pm\delta$. $\delta$
sets where ``outlier'' starts; about $1.35\sigma$ gives 95\% of OLS efficiency under Gaussian noise. Deep learning's ``smooth
L1'' (in object detection) is the same idea; ``gradient clipping'' is its cousin.
\end{qa}

\begin{qa}{\deep Can a classification model be turned into a regression model, and vice versa?}
\oneline{Yes: regression can be framed as classification over bins (ordinal or binned targets), and a classifier's
probability output is itself a regression of $\E[y]$ for a 0/1 target; many algorithms (trees, $k$-NN, SVM, neural nets) have
both forms with only the loss or leaf rule changed.}
\why Examples: predicting salary band (0--5, 5--10, 10--20 lakh) as multiclass when only bands are needed or when the
distribution is multimodal; ordinal regression (cumulative logit) for star ratings 1--5; decision trees use mean in leaves
for regression and majority vote for classification; SVM uses hinge or $\varepsilon$-insensitive loss. Binning loses ordering
information unless you use an ordinal model.
\end{qa}

\begin{qa}{What is deviance in a GLM, and how is it used?}
\oneline{Deviance is twice the difference in log-likelihood between a perfect (saturated) model and your model; it generalises
the residual sum of squares, and differences in deviance compare nested models.}
\why For Gaussian, deviance $=$ SSE$/\sigma^2$. For Poisson, $D = 2\sum[y\log(y/\mu) - (y - \mu)]$. Null deviance (intercept only)
minus residual deviance measures what the features explain; ``pseudo-$R^2$'' $= 1 - D_{model}/D_{null}$. Two nested models are
compared by the deviance difference, roughly $\chi^2$ with degrees of freedom equal to the number of extra parameters.
\worked Null deviance 1{,}200, model deviance 800: pseudo-$R^2 = 0.33$.
\end{qa}

\begin{qa}{\deep How would you model ``time to fill a vacancy'' when 30\% of jobs in your data are still open?}
\oneline{With a survival model, because open jobs are right-censored: we know they took at least their current age, and
dropping them or treating their age as the outcome biases the estimate downward.}
\why Ordinary regression on closed jobs only selects the fast-filling ones; using the current age of open jobs as if it were
the fill time understates it. Survival analysis uses the full information: Kaplan--Meier curves to describe, Cox proportional
hazards ($h(t\mid x) = h_0(t)e^{\mathbf w^\top\mathbf x}$) to model, or accelerated failure time models. Gradient boosting has
survival objectives too (\code{survival:cox} in XGBoost).
\end{qa}

\begin{qa}{What is an offset in Poisson regression? Give an example.}
\oneline{A term added to the linear predictor with its coefficient fixed at 1, typically $\log(\text{exposure})$, so the model
predicts a rate per unit of exposure.}
\why $\log\mu_i = \log t_i + \mathbf w^\top\mathbf x_i$ means $\mu_i = t_i e^{\mathbf w^\top\mathbf x_i}$: expected count is
exposure times rate. Example: listing enquiries on 99acres where listings were live 2 to 60 days: offset $\log(\text{days live})$.
Without it, ``old listings get more enquiries'' masquerades as a feature effect.
\end{qa}

\begin{qa}{\deep A colleague fits OLS to predict ``probability of hire'' and gets some predictions of 1.2. Fix the model
without changing to a black box.}
\oneline{Use a binomial GLM (logistic regression) so predictions stay in $(0,1)$; if they need a linear-probability reading, report
average marginal effects from the logistic model.}
\why The identity link has no range control; the logit link does. Average marginal effect $= \frac1n\sum w_j p_i(1-p_i)$ gives
``change in probability per unit of $x_j$'' averaged over the data, which is what the colleague wanted from OLS coefficients.
\end{qa}

\begin{qa}{When would you use isotonic regression?}
\oneline{When you know the relationship is monotone but not its shape, most commonly to calibrate classifier scores into
probabilities.}
\why Isotonic regression finds the best non-decreasing step function (pool adjacent violators algorithm). To calibrate a
gradient-boosted model, fit isotonic regression from its raw score to the 0/1 labels on a held-out set. It is flexible but can
overfit with little data; Platt scaling (a 1-feature logistic regression) is the low-data alternative.
\end{qa}

\begin{qa}{\deep Kernel ridge regression, SVR and Gaussian process regression all use kernels. How do they differ?}
\oneline{They share the kernel function but differ in loss and output: kernel ridge uses squared loss with a closed-form solution,
SVR uses the $\varepsilon$-insensitive loss giving a sparse solution, and a GP is the Bayesian version of kernel ridge that also
returns predictive uncertainty.}
\why Kernel ridge: $\hat f(\mathbf x) = \sum_i\alpha_iK(\mathbf x_i,\mathbf x)$ with $\boldsymbol\alpha = (K + \lambda I)^{-1}\mathbf y$;
every training point has non-zero $\alpha_i$. SVR: most $\alpha_i = 0$ (points in the tube), so prediction is cheaper and small
errors are ignored. GP: the posterior mean equals kernel ridge with $\lambda = $ noise variance; the posterior variance tells you
where the model is unsure (useful for active learning or pricing risk). All scale as $O(n^3)$ naively.
\end{qa}

\begin{qa}{How does XGBoost or LightGBM do Poisson, Gamma or quantile regression?}
\oneline{By swapping the loss: the trees fit gradients and Hessians of the chosen negative log-likelihood (Poisson, Gamma, Tweedie)
or of the pinball loss, usually on the log scale, so the boosting machinery is unchanged.}
\why For Poisson with log link and raw score $F$, the loss per row is $e^F - yF$; gradient $e^F - y = \mu - y$; Hessian $e^F =
\mu$. These plug straight into the XGBoost split-gain formula (Chapter 8). Parameters: \code{objective="count:poisson"},
\code{"reg:gamma"}, \code{"reg:tweedie"}, \code{"reg:quantileerror"} with \code{quantile\_alpha}.
\end{qa}
